\section{Tính toán và phân tích chỉ số hiệu suất}

    \begin{enumerate}
        \renewcommand{\labelenumi}{\arabic{enumi}.}

        \item \textbf{Kiến trúc hệ thống và chuỗi truyền dẫn dữ liệu}

        Hệ thống được thiết kế theo kiến trúc Dual-MCU nhằm tối ưu hóa việc xử lý song song và tăng độ tin cậy. Luồng dữ liệu di chuyển qua một ``đường ống'' (\textit{pipeline}) gồm bốn chặng nối tiếp:

        \begin{itemize}
            \item \textbf{Chặng 1 -- Air Interface:} Node ngoại vi phát dữ liệu qua sóng vô tuyến (Zigbee, LoRa, BLE) đến module thu tại Gateway.
            \item \textbf{Chặng 2 -- Module to MCU:} Module vô tuyến chuyển tiếp khung dữ liệu sang LAN-MCU (ESP32-S3) qua giao tiếp UART.
            \item \textbf{Chặng 3 -- Inter-MCU:} LAN-MCU gom dữ liệu từ các giao thức theo chu kỳ batch rồi chuyển toàn bộ sang WAN-MCU (ESP32-S3) qua bus SPI tốc độ cao.
            \item \textbf{Chặng 4 -- Cloud Interface:} WAN-MCU tuần tự hóa dữ liệu sang định dạng JSON và publish lên MQTT Broker qua kết nối Wi-Fi.
        \end{itemize}

        Mỗi chặng đều có độ trễ và giới hạn băng thông riêng. Nút thắt (bottleneck) của toàn hệ thống là chặng có thông lượng thấp nhất; phần tiếp theo sẽ định lượng từng chặng.

        % ---------------------------------------------------------------
        \item \textbf{Thông lượng dữ liệu và hiệu suất băng thông ($\eta$)}

        Thông lượng ứng dụng thực tế $R_{app}$ là lượng dữ liệu hữu ích (payload) còn lại sau khi trừ overhead giao thức tại tất cả các lớp. Hiệu suất băng thông được định nghĩa:
        \[
            \eta = \frac{R_{app}}{R_{PHY}}
        \]

        \textbf{Zigbee (E18-ZG120 / CC2530).} Tốc độ vật lý IEEE 802.15.4 tại băng 2,4~GHz đạt $R_{PHY} = 250\,\text{kbps}$ \cite{ieee802154_2020}. Mỗi frame Zigbee mang tối đa 99 byte payload ứng dụng sau khi trừ 28 byte overhead của các lớp MAC/NWK/APS trên tổng PSDU 127 byte \cite{zigbee_spec_r21}. Một chu kỳ phát--nhận ACK bao gồm: thời gian truyền frame ($127\,\text{B} \times 8 / 250\,\text{kbps} = 4{,}064\,\text{ms}$), khoảng LIFS ($0{,}192\,\text{ms}$) và frame ACK ($0{,}544\,\text{ms}$), tổng $T_{air} \approx 4{,}8\,\text{ms}$.

        Nút thắt thực tế không phải là air interface mà là giao tiếp UART từ module E18-ZG120 lên LAN-MCU. Module đóng gói mỗi khung Zigbee (127 byte PSDU) vào khung HEX của Ebyte với 19 byte framing overhead, tạo ra frame UART 146 byte. Tại tốc độ $115200\,\text{bps}$ với 10 bit/byte (8N1), thời gian truyền một frame là:
        \[
            T_{UART,Zigbee} = \frac{146 \times 10}{115200} \approx 12{,}67\,\text{ms}
        \]
        Điều này khiến UART chậm hơn air interface $12{,}67\,\text{ms} / 4{,}8\,\text{ms} \approx 2{,}6\times$, và trở thành điểm nghẽn của toàn chặng Zigbee.

        Trong kiểm thử thực tế với lệnh \texttt{reportTemperature()} -- chỉ mang 2--4 byte dữ liệu hữu ích nhưng phải gánh toàn bộ overhead ZCL -- thông lượng ứng dụng chỉ đạt $R_{app} \approx 8{,}2\,\text{kbps}$. Khi nhiều node phát đồng thời, cơ chế CSMA/CA ép các node lùi thời gian phát ngẫu nhiên (random backoff) khiến tổng thông lượng không tăng mà còn bị bão hòa tại ngưỡng này; vượt qua đó hàng đợi APS sẽ tràn và xảy ra mất gói (packet loss).

        \textbf{LoRa (Wio-E5 / STM32WLE5JC) -- Chế độ TEST/P2P.} Hệ thống sử dụng chế độ \texttt{AT+MODE=TEST} thay vì LoRaWAN, cho phép bỏ qua hoàn toàn giới hạn duty cycle 1\% của ETSI và stack LoRaWAN; thông lượng chỉ bị giới hạn bởi Time-on-Air. Tốc độ PHY tại SF7/BW125\,kHz đạt $R_{PHY} = 5{,}47\,\text{kbps}$ \cite{semtech_lora_design_guide}. Với gói 50 byte, Time-on-Air tính theo công thức Semtech là $T_{air} \approx 97{,}5\,\text{ms}$. Mỗi chu kỳ phát còn bao gồm thời gian xử lý AT command và nhận phản hồi \texttt{+TEST: TX DONE} ($\approx 16{,}1\,\text{ms}$), tổng $T_{cycle} \approx 113{,}6\,\text{ms}$, dẫn đến:
        \[
            R_{app,LoRa} = \frac{50 \times 8}{113{,}6\,\text{ms}} \approx 3{,}52\,\text{kbps}
        \]
        Giới hạn payload tối đa qua AT command là 169 byte/gói do độ dài lệnh tối đa 528 ký tự \cite{seeed_wioe5_at_cmd}.

        \textbf{BLE (Native ESP32-S3).} BLE được triển khai trực tiếp trên LAN-MCU nên không có tầng UART trung gian. Thông lượng bị giới hạn bởi Connection Interval ($CI = 20\,\text{ms}$) và kích thước ATT payload. Ở cấu hình mặc định ATT MTU = 23 byte (20 byte dữ liệu), thông lượng chỉ đạt $\approx 8\,\text{kbps}$ \cite{bluetooth_core_v5}. Khi kích hoạt Data Length Extension (DLE) và đàm phán MTU lên 247 byte (244 byte ATT payload), thông lượng cải thiện đáng kể:
        \[
            R_{app,BLE,DLE} = \frac{244 \times 8}{20\,\text{ms}} = 97{,}6\,\text{kbps}
        \]

        Bảng~\ref{tab:throughput_efficiency} tổng hợp so sánh thông lượng và hiệu suất của ba giao thức.

        \begin{longtable}{|p{3.2cm}|p{2.3cm}|p{2.7cm}|p{1.4cm}|p{4.2cm}|}
            \caption{So sánh thông lượng và hiệu suất của các giao thức.}
            \label{tab:throughput_efficiency} \\
            \hline
            \textbf{Giao thức} & $\boldsymbol{R_{PHY}}$ & $\boldsymbol{R_{app}}$ \textbf{thực tế} & $\boldsymbol{\eta}$ & \textbf{Nút thắt chính} \\ \hline
            \endfirsthead
            \hline
            \textbf{Giao thức} & $\boldsymbol{R_{PHY}}$ & $\boldsymbol{R_{app}}$ \textbf{thực tế} & $\boldsymbol{\eta}$ & \textbf{Nút thắt chính} \\ \hline
            \endhead
            \hline
            \endfoot
            \endlastfoot
            Zigbee & $250\,\text{kbps}$ & $\approx 8{,}2\,\text{kbps}$ & $3{,}3\%$ & CSMA/CA contention \& UART 12,67\,ms/frame \\ \hline
            LoRa SF7/125 (P2P) & $5{,}47\,\text{kbps}$ & $\approx 3{,}5\,\text{kbps}$ & $64\%$ & Time-on-Air ($\approx 97{,}5\,\text{ms/gói}$) \\ \hline
            BLE (không DLE) & $1000\,\text{kbps}$ & $\approx 8\,\text{kbps}$ & $0{,}8\%$ & MTU = 20 byte, $CI = 20\,\text{ms}$ \\ \hline
            BLE (bật DLE) & $1000\,\text{kbps}$ & $\approx 97{,}6\,\text{kbps}$ & $9{,}8\%$ & $CI = 20\,\text{ms}$ \\ \hline
        \end{longtable}

        % ---------------------------------------------------------------
        \item \textbf{Độ trễ End-to-End}

        Độ trễ End-to-End $T_{e2e}$ là tổng thời gian từ khi node gửi dữ liệu đến khi bản tin xuất hiện trên MQTT Broker:
        \[
            T_{e2e} = T_{air} + T_{UART} + T_{CI} + T_{batch} + T_{SPI} + T_{network}
        \]

        Các thành phần được định lượng như sau:

        \begin{itemize}
            \item \textbf{$T_{air}$ -- Thời gian truyền sóng:} Zigbee $\approx 4{,}8\,\text{ms}$ (127 byte frame + LIFS + ACK), LoRa SF7 $\approx 97{,}5\,\text{ms}$ (ToA, gói 50 byte), BLE $\approx 2{,}2\,\text{ms}$ (LL PDU 263 byte tại 1 Mbps + turnaround).
            \item \textbf{$T_{UART}$ -- Độ trễ giao tiếp module--MCU:} Zigbee $\approx 12{,}67\,\text{ms}$ (146 byte Ebyte frame tại 115200 bps, 8N1), LoRa $\approx 16{,}1\,\text{ms}$ (AT command + phản hồi TX DONE). BLE không có thành phần này do tích hợp native trên LAN-MCU.
            \item \textbf{$T_{CI}$ -- Độ trễ chờ Connection Interval (chỉ BLE):} Vì BLE chỉ truyền dữ liệu tại các thời điểm cố định cách nhau $CI = 20\,\text{ms}$, một gói đến bất kỳ thời điểm nào trong chu kỳ phải chờ đến interval tiếp theo. Độ trễ chờ này dao động $0$--$CI = 0$--$20\,\text{ms}$, trung bình $CI/2 = 10\,\text{ms}$. Zigbee và LoRa không có thành phần này.
            \item \textbf{$T_{batch}$ -- Độ trễ gom gói:} Dao động $0$--$10\,\text{ms}$ tùy thời điểm dữ liệu đến trong chu kỳ batch flush của LAN-MCU.
            \item \textbf{$T_{SPI}$ -- Độ trễ bus Inter-MCU:} Tại xung nhịp 60\,MHz với DMA, worst-case burst ($\approx 5193$ byte, xem Mục~4) mất $\approx 0{,}693\,\text{ms}$; hệ số sử dụng bus chỉ đạt $6{,}93\%$.
            \item \textbf{$T_{network}$ -- Độ trễ WAN lên Cloud:} Trung bình $\approx 30\,\text{ms}$ qua Wi-Fi đến MQTT Broker.
        \end{itemize}

        Với BLE, độ trễ worst-case cần tính thêm $T_{CI,max} = 20\,\text{ms}$ (gói đến ngay sau khi vừa bỏ lỡ một interval):
        \[
            T_{e2e,BLE,worst} = T_{air} + T_{CI,max} + T_{batch,max} + T_{SPI} + T_{network}
            \approx 2{,}2 + 20 + 10 + 0{,}7 + 30 \approx 63\,\text{ms}
        \]

        \begin{longtable}{|p{4.5cm}|p{4.5cm}|p{5cm}|}
            \caption{Độ trễ dự kiến theo từng kịch bản truyền.}
            \label{tab:latency_estimation} \\
            \hline
            \textbf{Kịch bản truyền} & \textbf{Độ trễ dự kiến} & \textbf{Thành phần chiếm ưu thế} \\ \hline
            \endfirsthead
            \hline
            \textbf{Kịch bản truyền} & \textbf{Độ trễ dự kiến} & \textbf{Thành phần chiếm ưu thế} \\ \hline
            \endhead
            \hline
            \endfoot
            \endlastfoot
            Zigbee $\rightarrow$ MQTT & $48\,\text{ms}$--$58\,\text{ms}$ & $T_{UART}$ (12,67 ms) \& $T_{network}$ (30 ms) \\ \hline
            LoRa SF7 (P2P) $\rightarrow$ MQTT & $144\,\text{ms}$--$154\,\text{ms}$ & $T_{air}$ (97,5 ms) -- chiếm $\approx 65\%$ \\ \hline
            BLE $\rightarrow$ MQTT (trung bình) & $33\,\text{ms}$--$43\,\text{ms}$ & $T_{network}$ (30 ms) -- chiếm $\approx 80\%$ \\ \hline
            BLE $\rightarrow$ MQTT (worst-case) & $53\,\text{ms}$--$63\,\text{ms}$ & $T_{CI,max}$ (20 ms) \& $T_{network}$ (30 ms) \\ \hline
        \end{longtable}

        % ---------------------------------------------------------------
        \item \textbf{Khả năng chịu tải và phân tích bus SPI}

        Để đánh giá mức độ sử dụng bus SPI, ta tính lượng dữ liệu tích lũy tại LAN-MCU trong một batch interval $T_{batch} = 10\,\text{ms}$ dựa trên thông lượng thực tế của từng giao thức (Bảng~\ref{tab:batch_data}).

        \begin{longtable}{|p{3cm}|p{4cm}|p{6.5cm}|}
            \caption{Dữ liệu tích lũy theo giao thức trong một batch interval.}
            \label{tab:batch_data} \\
            \hline
            \textbf{Giao thức} & \textbf{Thông lượng} & \textbf{Dữ liệu/batch ($T_{batch} = 10\,\text{ms}$)} \\ \hline
            \endfirsthead
            \hline
            \textbf{Giao thức} & \textbf{Thông lượng} & \textbf{Dữ liệu/batch ($T_{batch} = 10\,\text{ms}$)} \\ \hline
            \endhead
            \hline
            \endfoot
            \endlastfoot
            Zigbee     & $8{,}2\,\text{kbps}$   & $8200 / 8 \times 0{,}010 \approx 10\,\text{byte}$   \\ \hline
            BLE (DLE)  & $97{,}6\,\text{kbps}$  & $97600 / 8 \times 0{,}010 \approx 122\,\text{byte}$ \\ \hline
            LoRa SF7/125 & $3{,}5\,\text{kbps}$ & $3500 / 8 \times 0{,}010 \approx 4\,\text{byte}$    \\ \hline
            \textbf{Tổng} & $\mathbf{109{,}3\,\text{kbps}}$ & $\mathbf{\approx 137\,\text{byte}}$ \\ \hline
        \end{longtable}

        Với 137 byte mỗi batch, thời gian SPI cần thiết là:
        \[
            T_{SPI,batch} = \frac{137 \times 8}{60 \times 10^6} \approx 0{,}0183\,\text{ms}
        \]

        Hệ số sử dụng bus SPI trong điều kiện trung bình:
        \[
            U_{SPI} = \frac{T_{SPI,batch}}{T_{batch}} = \frac{0{,}0183}{10} \approx 0{,}183\%
        \]

        Để đánh giá điều kiện worst-case, ta giả thiết tất cả node gửi đồng thời với frame kích thước tối đa. Hệ thống hỗ trợ tối đa 20 node Zigbee, 8 node BLE, và 1 node LoRa (xem ràng buộc phần sau). Burst tổng được tính:

        \begin{longtable}{|p{3cm}|p{2cm}|p{3.5cm}|p{4.5cm}|}
            \caption{Tính toán burst SPI trong điều kiện worst-case.}
            \label{tab:spi_worst} \\
            \hline
            \textbf{Giao thức} & \textbf{Số node} & \textbf{Frame tối đa/node} & \textbf{Tổng burst} \\ \hline
            \endfirsthead
            \hline
            \textbf{Giao thức} & \textbf{Số node} & \textbf{Frame tối đa/node} & \textbf{Tổng burst} \\ \hline
            \endhead
            \hline
            \endfoot
            \endlastfoot
            Zigbee   & 20 & 146 byte (PSDU + Ebyte overhead) & $20 \times 146 = 2920\,\text{byte}$ \\ \hline
            BLE (DLE) & 8  & 263 byte (LL PDU tối đa)         & $8 \times 263 = 2104\,\text{byte}$  \\ \hline
            LoRa P2P & 1  & 169 byte (AT command limit)       & $1 \times 169 = 169\,\text{byte}$   \\ \hline
            \textbf{Tổng} & \textbf{29} & -- & $\mathbf{5193\,\text{byte}}$ \\ \hline
        \end{longtable}

        Thời gian SPI cần thiết và hệ số sử dụng trong điều kiện worst-case:
        \[
            T_{SPI,worst} = \frac{5193 \times 8}{60 \times 10^6} \approx 0{,}693\,\text{ms}
        \]
        \[
            U_{SPI,worst} = \frac{0{,}693}{10} \approx 6{,}93\%
        \]

        Thông lượng tối đa lý thuyết của SPI là:
        \[
            R_{SPI,max} = 60\,\text{Mbps} = 7{,}5\,\text{MB/s}
        \]

        Ngoài giới hạn bus, mỗi giao thức còn chịu ràng buộc về số node:

        \begin{itemize}
            \item \textbf{Zigbee (E18-ZG120 / CC2530):} CC2530 có $8\,\text{KB}$ SRAM; sau khi Z-Stack chiếm $\approx 5\,\text{KB}$, phần còn lại chỉ đủ cho tối đa 20 node trực tiếp theo cấu hình mặc định (\texttt{NWK\_MAX\_DEVICE\_LIST = 20}). Vượt ngưỡng này sẽ phát sinh lỗi \texttt{0x18} (Not enough cache) hoặc \texttt{0x11} (Memory full).
            \item \textbf{LoRa (P2P):} Chế độ point-to-point chỉ hỗ trợ một liên kết tại một thời điểm; không có ràng buộc mạng lưới như LoRaWAN.
            \item \textbf{BLE:} NimBLE stack trên ESP32-S3 giới hạn tối đa 8 kết nối đồng thời.
        \end{itemize}

        \textbf{Kết luận:} Bus SPI hoàn toàn không phải nút thắt -- hệ số sử dụng trung bình chỉ đạt $0{,}183\%$, worst-case $6{,}93\%$ (khi toàn bộ 29 node gửi frame tối đa đồng thời). Nguy cơ suy giảm PDR xuất phát từ Zigbee Listener buffer ($512\,\text{byte}$) bị tràn khi LAN-MCU bận phục vụ ngắt SPI/UART tần suất cao (mỗi $10\,\text{ms}$); đây là điểm cần tối ưu để đạt mục tiêu PDR toàn hệ thống $> 85\%$ dưới tải nặng.

    \end{enumerate}